\section{Whole expression centers in eleven integers}
\label{sec:scalar}
\subsection{Complete polynomial definitions}
The independent coordinates are exactly
\[
 (j,C,z,h,r,v,w,T,c,n,s)\in\Z^{11}.
\]
The six fixed ports are the authentic source34 numerals specified in
Section~\ref{sec:foundation}; none is another variable. Define the following
ordinary polynomial abbreviations:
\begin{align*}
 J&=6j+1,& q&=\beta J+1,& W&=12w+4,\\
 X&=qW,&Y&=v,&E&=XY,\\
 a&=E+Y,&A&=a+2,&H&=4a+3,\\
 \Delta&=a^2+H=A^2-1,&&&
\end{align*}
\begin{align*}
 F&=(K_0+W)C-z(q-1),\\
 R&=[q(q-F)-C](q^2-1)+(M_C+qM_F)J,\\
 k&=R+1+hE,\qquad \eta=c-Yk,\\
 D_{\rm err}&=(\mathrm{cone\_A}F-q)^2+(\mathrm{cone\_B}C-q)^2,\\
 B_0&=k^2(q^2-D_{\rm err})-4q^2\eta^2,\\
 \mathcal B&=Y^2B_0-k^2q^8,\qquad D=\mathcal B-4,\\
 \ell&=XY^2,\qquad D_t=1+\ell(\ell+1)k^2,\\
 D_g&=1+DD_ts^2,\qquad P=D_tD_g.
\end{align*}
The cofactor, auxiliary norm and carrier registers are
\begin{align*}
 D_m&=1+\Delta c^2,& B&=rAD_m,& U&=rc^2D_m,\\
 I_2&=c^4B=Ac^2U,& D_f&=1+\Delta I_2,& Q&=\Delta^2I_2,\\
 \delta&=Q-1,& V&=D_f(c-R+4AUT)-c,& C_y&=QV^2-1,\\
 K&=X(3c-2X)+2,&M&=2q^3,&D_c&=HM,\\
 S&=srHM^5,&&&
\end{align*}
\[
 z_c=D_cT+MK+H(v-q^3),\qquad
 A_c=(z_c^2+3D_c^2)(z_c^3-19D_c^3).
\]
Choose one fixed even weight $L$ and one center $\Psi$ below:
\begin{align*}
 \Psi_{316}&=[(8\Delta+4)U-rc^4]V,&L_{316}&=8,\\
 \Psi_{308}&=\Psi_{316}+2U^3T(c^2-2),&L_{308}&=8,\\
 \Psi_{304}&=\Psi_{316}+U^2T[2U(c^2-2)+16A^3],&L_{304}&=8,\\
 \Psi_{300}&=[(32\Delta+16)U-4rc^4]V&&\\
 &\quad+U^2T[8U(c^2-2)+64A^3+3rc^6],&L_{300}&=32.
\end{align*}
Every center is an integer polynomial independent of $n$. In each case put
\begin{align}
 Z&=S(n+\Psi)-A_c,\nonumber\\
 p_a&=\delta Z^2,\qquad p_b=\delta S^2P,
 \qquad p_c=L^2S^2BD_fC_y,\nonumber\\
 \mathcal N_{L,\Psi}&=(p_a+p_b-p_c)^2-4p_ap_b.
 \label{eq:scalarfinal}
\end{align}
The equation is $\mathcal N_{L,\Psi}=0$ over $\Z^{11}$. No division,
radical or variable exponent is an operation of this polynomial. The
literal graph shares the identities
\[
 P=D_t+(\mathcal B-4)(D_ts)^2,
\]
\[
 \mathcal B=(Yk)^2(q^2-D_{\rm err})
       -q^2\bigl[4(Y\eta)^2+(kq^3)^2\bigr],
 \qquad A^2=\Delta+1.
\]
Sharing these registers changes graph cost, not the polynomial.

\subsection{An even weight and integral center lemma}
\begin{theorem}\label{thm:evenweight}
For every fixed nonzero even $L$ and every integer polynomial $\Psi$ in
the ten coordinates other than $n$, the projection of the integer zero
set of \eqref{eq:scalarfinal} onto those ten coordinates equals that of
the retained Round~17 free-scale equation. The latter uses
$L=2$ and $\Psi=(2\Delta+1)UV$.
\end{theorem}
\begin{proof}
We first record the unconditional arithmetic facts from Sections 2--4
of the retained \file{LINEAR_COFACTOR_ROUND17_PROOF.md}. They do not
assume a zero of the present norm. The authentic affine restrictions give
\begin{gather*}
 q\equiv0\pmod2,\qquad q\equiv2\pmod3,\qquad |q|\ge16,\\
 X\equiv0\pmod8,\qquad X\equiv2\pmod3,\qquad |X|\ge4|q|.
\end{gather*}
Thus $A\ne0$, $\Delta=A^2-1>1$, $\Delta\equiv0\pmod3$,
$H$ is odd and nonzero, and $D_m,D_t\ge1$. The exact modulo-sixteen
budget identity gives $D\equiv8$ or $12\pmod{16}$, so $D$ is nonzero
and nonsquare. Also $P\equiv1\pmod4$ and $P\ne0$.
Both $D_f$ and $\delta$ are nonzero, $Q$ is even, and $\delta$ is odd.

The rational polynomial $\pi(t)=(t^2+3)(t^3-19)$ has no rational root.
Since $D_c\ne0$,
\[
 A_c=D_c^5\pi(z_c/D_c)\ne0.
\]
If $S=0$, the output norm is $(\delta A_c^2)^2\ne0$.
Every new zero therefore has $S,s,r,B\ne0$.

Put $D_a=\delta C_y$ and $C_*=BD_fD_a$. One has $D_a>0$:
if $c=0$, then $Q=0$ and $D_a=1$; if $c\ne0$ and $B<0$,
then $\delta,C_y<0$; if $c\ne0$ and $B>0$, then
$D_f>|c|+1$, so $V=D_fb-c$ is nonzero for every integer $b$, and
$\delta,C_y>0$. Moreover $BD_f>0$ whenever $c\ne0$.
Hence $C_*\ne0$, and its only possible negative case is $c=0,B<0$.

The four-conjugate norm identity for \eqref{eq:scalarfinal}, after
normalizing by nonzero $S$ and multiplying by $\delta$, gives a rational
signed sum
\begin{equation}\label{eq:radicalsum}
 \pm L\sqrt{C_*}\ \pm\delta\sqrt P\in\Q.
\end{equation}
If exactly one radicand is negative, its nonzero imaginary contribution
cannot cancel. If both are negative, then $c=0$, $\delta=-1$ and
$C_*=B$. Cancellation would force $L^2B=P$, impossible modulo four
because $L$ is even and $P\equiv1\pmod4$. Therefore $C_*,P>0$.

Write these positive integers as squares times positive squarefree integers.
Different nontrivial squarefree radicals are independent over $\Q$.
The only remaining possible irrational cancellation would have the same
nontrivial squarefree part in both terms of \eqref{eq:radicalsum}.
Its first coefficient would be even, whereas its second coefficient would
be odd: $L$ is even and $\delta,P$ are odd. Cancellation is impossible.
Thus $C_*$ and $P$ are positive integer squares.

The cofactor coprimalities are unchanged:
\[
 \gcd(B,D_f)=\gcd(B,D_a)=\gcd(D_f,D_a)=1.
\]
Indeed $D_f\equiv1\pmod B$ and $Q\equiv0\pmod B$, giving
$D_a\equiv1\pmod B$. Modulo $D_f$, one has
\[
 Q\equiv-\Delta,\quad\delta\equiv-A^2,\quad
 V\equiv-c,\quad C_y\equiv-D_m,\quad
 D_a\equiv A^2D_m.
\]
Since $B=rAD_m$, while $D_f\equiv1$ modulo both $A$ and $D_m$,
these imply the last coprimality. Therefore $|B|,|D_f|,D_a$ are squares.
As $D_a>0$, $B,D_f$ have the same sign. But $D_f\equiv1\pmod3$
excludes a negative square, so there are positive integers $u,f,e$ with
\begin{equation}\label{eq:scalarroots}
 B=u^2,\qquad D_f=f^2,\qquad D_a=e^2.
\end{equation}
Similarly $\gcd(D_t,D_g)=1$ and $D_t\ge1$ give
$D_t=\tau^2$, $D_g=g^2$ with $\tau,g>0$.
Since $s\ne0$, the guard forces $\mathcal B\ge4$, strengthened by its
modulo-sixteen class to $\mathcal B\ge12$.
It follows that $k,Y\ne0$, $B_0>0$, and
\[
 Y^2B_0>k^2q^8,\qquad B_0\le k^2q^2,
 \qquad |Y|>|q|^3.
\]
The centered strip from $B_0>0$ then excludes $c=0$.

Define proof-only integers $I=c^2u$ and $T_{\rm old}=uT$.
By \eqref{eq:scalarroots},
\[
 f^2-\Delta I^2=1,\qquad Q=(\Delta I)^2,\qquad
 V=f^2(c-R+4IT_{\rm old})-c.
\]
The rational number $y=e/\delta$ satisfies
\begin{equation}\label{eq:auxnorm}
 QV^2-(Q-1)y^2=1.
\end{equation}
It is integral by adjacent-coefficient descent. If $b$ is its reduced
denominator, then $b^2\mid Q-1$. After choosing signs, for $|V|>1$ the
unimodular transformation
\[
 V'=(2Q-1)V-2(Q-1)y,\qquad
 y'=(2Q-1)y-2QV
\]
preserves \eqref{eq:auxnorm} and the reduced denominator, with
$0<|V'|<|V|$. The first coordinate is integral because $(Q-1)y$ is
integral; denominator preservation also follows from
$\gcd(2Q-1,b)=1$. For $V>1$ and $Q>1$, the norm gives
\[
 V<y<\frac{(2Q-1)V}{2(Q-1)}.
\]
The squared upper bound minus $y^2$, multiplied by $4(Q-1)^2$,
is $V^2+4(Q-1)>0$. Consequently $0<V'<V$.
Descent reaches $|V|=1$, where $|y|=1$, forcing $b=1$.
The formal $c=0$ case would directly give $y=-1$, although it is already
excluded here.

The two normalized roots of the original four-conjugate equation are now
integers $Lufy$ and $\tau g$. Hence $Z/S$ is integral. Since $\Psi$
is integral, $S\mid A_c$. No size or sign condition on $\Psi$ was used.
Keep the ten outer coordinates of the new zero and set
\[
 n_{\rm old}=A_c/S+2ufy+\tau g-(2\Delta+1)UV.
\]
This is an integer and makes one old conjugate factor vanish, producing an
actual Round~17 zero. Conversely an old zero supplies these same integral
roots and divisibility by the retained theorem. Keeping its ten outer
coordinates, set
\[
 n_{\rm new}=A_c/S+Lufy+\tau g-\Psi.
\]
One new factor vanishes. These constructions prove equality of projected
zero sets. They need not be mutually inverse on $n$, since the norm can
have several sign branches. For a fixed $L$, changing only the center is
separately a literal bijective integral shear of $n$.
\end{proof}

\subsection{The complete inherited calibration chain}
Theorem~\ref{thm:evenweight} reduces both directions to an actual old
zero, not merely a list of necessary Pell properties. Thus its inherited
scale, sign, rank and syntactic interfaces remain exactly applicable.
For clarity, the retained source proof proceeds as follows.

Carrier divisibility first gives $M\mid H(v-q^3)$, even $Y$, and
\[
 H\mid q^{15}K_q^5,\qquad K_q=X(X+2ac)-1.
\]
The strong and auxiliary pairs, together with $c^2\mid I$, recover the
main Pell root by absolute folding. With $D_t=\tau^2$, the centered
strip, $|Y|>|q|^3$ and the authentic packing inequalities
\[
 |q|(|q|^2-1)<R<|q|^4-2,
\]
the signed rank theorem yields
\[
 c=\sgn(c)\psi_A(R),\qquad R=2t+1,
 \qquad |H|>|X|(|X|-3)^t.
\]
The powered projection at the actual large base forces
$\sgn(c)X=2^R$. Reduction modulo three fixes $c>0$ and $X=2^R$;
the old strip fixes $k,Y>0$. Exact $H\mid K$ restores
$Y=(2v_{\rm old}+1)q^3$. The uniform half-binomial theorem then restores
the positive native interval, and fresh odd strong witnesses produce an
actual signed Q41 source zero. Only that restored zero is passed to the
retained source34/native/compiler/finite-refutation calibration.

Conversely, canonical source witnesses give $\mathcal B>64$.
After choosing the guard pair $s,g$, the old proof chooses one strong
pair $I,f$ with $4sAdc^2\mid I$, where $D_m=d^2$, and sets
\[
 u=I/c^2,\qquad t_{\rm strong}=I/(Adc^2),\qquad
 r=At_{\rm strong}^2.
\]
That pair is kept fixed. The supported-modulus theorem interpolates
$T_{\rm old}$ modulo $usr$ using the same pair, while an intersective
root of $\pi$ modulo $sr$ makes $A_c/S$ integral. There is no reselection
of the strong pair after a residue request. The present change replaces
only its final integer $n$ by $n_{\rm new}$ above and introduces no new
interpolation or choice-order assumption. Full inherited proofs and their
source identities are retained separately; they are not promoted here to
a newly checked formal-kernel endpoint.

\subsection{Exact degrees of the entire output polynomials}
The degree claims concern the authentic fixed coefficient tuple, not
generic ports, line samples or syntactic upper bounds. The shared register
degrees are
\begin{center}\small
\begin{tabular}{@{}ll@{}}\toprule
Registers & Exact degrees\\\midrule
$a,A;\ \Delta;\ R,k$ & $3;\ 6;\ 5$\\
$\mathcal B,D_t;\ P$ & $18;\ 56$\\
$D_m,B,U;\ D_f;\ Q,\delta$ & $8,12,11;\ 22;\ 28$\\
$V,C_y;\ S,A_c;\ \Psi,Z$ & $37,102;\ 20,35;\ 54,74$\\
$p_a,p_b;\ 4p_ap_b$ & $176,124;\ 300$\\\bottomrule
\end{tabular}
\end{center}
The highest form of $k$ is that of $R$, since $hXY$ has degree four.
Put $W_0=Sn-A_c$, so $\deg W_0=35$. In the free polynomial ring
with independent coordinates $A,c,r,T,R$, assign weights $(3,1,1,1,5)$
and use the exact definitions $\Delta=A^2-1$, $D_m=1+\Delta c^2$,
$B=rAD_m$, $U=rc^2D_m$, $D_f=1+\Delta c^4B$ and
$V=D_f(c-R+4AUT)-c$. Define
\begin{align}
 D_{\rm center}&=\Psi^2-L^2BD_fV^2,\nonumber\\
 E_{\rm center}&=\delta\Psi^2-L^2BD_fC_y\nonumber\\
 &=\delta D_{\rm center}-L^2BD_f(V^2-1).
 \label{eq:centeridentity}
\end{align}
Exact expansion gives
\begin{center}\small
\begin{tabular}{@{}crlr@{}}\toprule
Center & $\deg D_{\rm center}$ & Highest weighted form & $\deg E_{\rm center}$\\\midrule
316&90&$-128A^{18}c^{28}r^6T^2$&118\\
308&86&$-1024A^{19}c^{22}r^5T^2$&114\\
304&84&$-48A^{16}c^{28}r^6T^2$&112\\
300&82&$4096A^{16}c^{26}r^6T^2$&110\\\bottomrule
\end{tabular}
\end{center}
Every omitted monomial has strictly smaller weight. The following finite
identities give a direct route to those expansions. Put
\[
 \lambda=(8\Delta+4)U-rc^4,\quad \Psi_0=\lambda V,\quad
 h_*=2U^3T(c^2-2),\quad b_*=16A^3U^2T,\quad e_*=3rc^6U^2T.
\]
Then
\begin{align*}
 D_{316}&=[8Ur(2c^2-c^4)+r^2c^8-64B]V^2,\\
 D_{308}&=D_{316}+2\Psi_0h_*+h_*^2,\\
 D_{304}&=D_{308}+2(\Psi_0+h_*)b_*+b_*^2,\\
 D_{300}&=16D_{304}+8(\Psi_0+h_*+b_*)e_*+e_*^2.
\end{align*}
They follow by squaring the four centers and using
$BD_f=B+\Delta A^2U^2$. Substitution of
\[
 U=r(c^4A^2+c^2-c^4)
\]
and the displayed formula for $V$ gives the complete finite expansion;
the independent checker supplies every exact remainder coefficient.
The second term of \eqref{eq:centeridentity} has weight 108, smaller than
each leading $\delta D_{\rm center}$. Since the highest form of
$\delta$ is $rA^7c^6$, the respective highest forms of $E_{\rm center}$ are
\[
 -128A^{25}c^{34}r^7T^2,\quad
 -1024A^{26}c^{28}r^6T^2,\quad
 -48A^{23}c^{34}r^7T^2,\quad
 4096A^{23}c^{32}r^7T^2.
\]

Now the exact norm cut is
\begin{equation}\label{eq:normcut}
 p_a+p_b-p_c=\delta W_0^2+2\delta S\Psi W_0+
             \delta S^2P+S^2E_{\rm center}.
\end{equation}
The first three terms have degrees at most 98, 137 and 124.
The last has exact degree 158, 154, 152 or 150, respectively, with a
unique highest monomial after the actual substitution $A=a+2$.
Squaring therefore gives degrees 316, 308, 304 and 300. In the first
three cases, $4p_ap_b$ has smaller degree 300. In the last case its highest
form has $T$-degree two, whereas the highest monomial of the squared inner
term has $T$-degree four. That monomial cannot cancel.

For explicit coefficients, put $q_0=6\beta j$ and $a_0=12q_0wv$.
Then $S_0=128sra_0q_0^{15}$. Surviving highest monomials have the following
exponents, with all omitted variable exponents zero:
\begin{center}\small
\begin{tabular}{@{}cll@{}}\toprule
Degree & Exponents in $(j,w,v,s,r,c,T)$ & Positive coefficient\\\midrule
316&$(114,54,54,4,18,68,4)$&$2^{264}3^{168}\beta^{114}$\\
308&$(116,56,56,4,16,56,4)$&$2^{276}3^{172}\beta^{116}$\\
304&$(110,50,50,4,18,68,4)$&$2^{246}3^{162}\beta^{110}$\\
300&$(110,50,50,4,18,64,4)$&$2^{262}3^{160}\beta^{110}$\\\bottomrule
\end{tabular}
\end{center}
Each exponent sum is the stated degree and each coefficient is nonzero
for the actual $\beta$. For degree 300, the uncancelled $T^4$ summand before
expanding $a_0,q_0$ is
\[
 2^{52}s^4r^{18}c^{64}T^4a_0^{50}q_0^{60}.
\]
This establishes exact attained degrees. It is not a lower bound for all
possible eleven-variable constructions.

\subsection{Live costs and finite checks}
The actual graphs have the following operation counts:
\begin{center}\small
\begin{tabular}{@{}crrrr@{}}\toprule
Degree & Additions & Subtractions & Multiplications & Total\\\midrule
316&27&21&85&133\\
308&28&22&90&140\\
304&30&22&92&144\\
300&31&22&95&148\\\bottomrule
\end{tabular}
\end{center}
A square or multiplication by a fixed numeral costs one binary gate;
register reuse is free. All eleven variables and six ports are live.
The independently implemented checker imports neither the producer nor
its graph class. It binds all named registers and live gates in free
polynomial rings, verifies the four-conjugate identity, center remainders,
degree cuts and actual leading coefficients, authenticates and reconstructs
all six numerical payloads, checks source pins and liveness, and rejects
altered constants, coordinates, degree claims and selected semantic gates.
These are finite algebraic checks, distinct from the retained unbounded
Pell, interpolation, intersectivity and native-refutation theorems.
